\chapter{Smoothed-aggregation multigrid}
\label{ch:amg}
\impl{src/gpu\_amg.jl, src/linear\_solvers.jl}

\section{Why multigrid}

Every preconditioner of Chapter~\ref{ch:preconditioners} is local: it couples
each unknown to its immediate algebraic neighbours.  Local operators damp
high-frequency error efficiently and low-frequency error not at all, which is
exactly why $\kappa(\Kmat)$ grows like $(L/h)^2$ under refinement --- the
smoothest error modes are the ones a local operator cannot see.

Multigrid removes what smoothing cannot by representing the smooth error on a
coarser problem, where it is no longer smooth relative to the grid.  The payoff
is an iteration count that is nearly \emph{independent} of problem size.  On
a $530$k-DOF torsion problem, algebraic multigrid holds $5$--$17$ conjugate
gradient iterations across a range of $\Delta t$ over which Jacobi grows from
$30$ to $442$.

Carina uses \emph{algebraic} multigrid: the hierarchy is built from the matrix
alone, with no geometric coarse mesh, which is what makes it applicable to
unstructured meshes.

\section{The near-nullspace}
\label{sec:nullspace}

Algebraic multigrid needs to know which error modes the smoother will not
remove, so that the coarse space can represent them.  For elasticity these are
the modes with no strain energy: the six rigid-body modes --- three
translations and three rotations.  They form the columns of the
\emph{near-nullspace} $\Bnull \in \mathbb{R}^{3n_{\mathrm{node}} \times 6}$,
built from nodal coordinates about the centroid $\bar{\vect{x}}$:
\begin{equation}
  \Bnull = \left[\
  \begin{array}{ccc|ccc}
    1 & 0 & 0 & 0     & \hat{z}  & -\hat{y} \\
    0 & 1 & 0 & -\hat{z} & 0     & \hat{x}  \\
    0 & 0 & 1 & \hat{y}  & -\hat{x} & 0
  \end{array}\ \right],
  \qquad \hat{\vect{x}} = \vect{x} - \bar{\vect{x}},
\end{equation}
one such block per node, restricted to the free degrees of freedom.

\paragraph{Current configuration, not reference.}  Carina builds $\Bnull$ from
the \emph{current} nodal coordinates $\vect{X} + \disp$ at every hierarchy
build, not from the reference configuration.  At finite deformation the true
near-nullspace of $\Kmat(\Uvec)$ consists of rigid-body motions about the
\emph{deformed} body; a nullspace frozen at the reference configuration
describes rotations the body no longer has.  The consequence is not subtle:
on a twisted-bar run, a frozen reference nullspace degraded the iteration count
from $63$ to over $400$ as the bar rotated, while the hierarchy build count
stayed at one and nothing in the log indicated a problem.

\section{Hierarchy construction}

Each level is built in four stages.

\paragraph{1. Strength of connection.}  A symmetric strength measure filters
$\Kmat$ into a graph $\tensor{S}$ retaining only algebraically strong
couplings, so that aggregation follows the directions in which error actually
propagates.

\paragraph{2. Aggregation.}  Standard aggregation partitions the nodes of
$\tensor{S}$ into disjoint aggregates, each becoming one coarse node.

\paragraph{3. Tentative prolongator.}  For each aggregate $g$, the restriction
of the near-nullspace to that aggregate's rows is factored,
\begin{equation}
  \Bnull|_g = \tensor{Q}_g \tensor{R}_g,
\end{equation}
by a thin QR.  The orthonormal $\tensor{Q}_g$ becomes the aggregate's block of
the tentative prolongator $\tensor{T}$, and $\tensor{R}_g$ becomes the coarse
near-nullspace $\Bnull_c$ for the next level.  This construction guarantees
\begin{equation}
  \Bnull = \tensor{T}\Bnull_c ,
\end{equation}
i.e.\ the coarse space represents every near-null mode \emph{exactly} --- which
is the whole point of supplying $\Bnull$.

\begin{remark}
An aggregate with fewer nodes than $\Bnull$ has columns is rank deficient, and
its QR yields fewer than six usable columns.  The implementation handles this
by giving such an aggregate $r = \min(\abs{g}, 6)$ columns and leaving the
remainder structurally empty, rather than producing near-zero columns that
would inflate every downstream sparse product.
\end{remark}

\paragraph{4. Prolongator smoothing.}  The tentative prolongator is improved by
one damped Jacobi sweep,
\begin{equation}
  \Pmat = \left(\Ident - \frac{4}{3\lambda_{\max}}\Dmat^{-1}\Kmat\right)\tensor{T},
\end{equation}
which lowers the energy of the coarse basis functions and is what distinguishes
\emph{smoothed} aggregation from plain aggregation.  Restriction is
$\Rrest = \Pmat^{\mathsf{T}}$, and the coarse operator is the Galerkin triple
product
\begin{equation}
  \Kmat_c = \Rrest\, \Kmat\, \Pmat .
  \label{eq:galerkin}
\end{equation}

\begin{remark}[Memory]
Evaluated naively, \eqref{eq:galerkin} preallocates from an up-front
nonzero-count estimate that is wildly pessimistic for a six-column
near-nullspace --- demanding tens of gigabytes at $1.5$M DOF for a product
whose true size is about $4$~GB.  Carina evaluates the fine-level product in
column slabs of $\Pmat$, so peak transient memory is one slab's product rather
than the whole estimate, and hands the much smaller coarse problem to the
standard routine for the remaining levels.
\end{remark}

\section{The V-cycle}

One V($\nu,\nu$)-cycle applied to a residual $\vect{r}$ produces a correction
$\vect{z}$:

\begin{algorithm}[h]
\caption{V-cycle, level $\ell$}
\begin{algorithmic}[1]
  \STATE $\vect{x} \gets \vect{0}$
  \STATE pre-smooth $\nu$ times on $\Kmat^{\ell}\vect{x} = \vect{b}^{\ell}$
  \STATE $\vect{b}^{\ell+1} \gets \Rrest^{\ell}\left(\vect{b}^{\ell} - \Kmat^{\ell}\vect{x}\right)$
  \STATE recurse to level $\ell+1$, or solve directly if coarsest
  \STATE $\vect{x} \gets \vect{x} + \Pmat^{\ell}\vect{x}^{\ell+1}$
  \STATE post-smooth $\nu$ times
\end{algorithmic}
\end{algorithm}

The smoother is damped Jacobi with $\omega = 4/(3\lambda_{\max})$, the same
weight used for prolongator smoothing, with $\lambda_{\max}$ estimated per
level by power iteration.  The coarsest level is solved by a dense
pseudo-inverse, formed once at build time; at that size a dense solve is
cheaper than another level of hierarchy.

Symmetry of the cycle --- equal pre- and post-smoothing, and
$\Rrest = \Pmat^{\mathsf{T}}$ --- is what makes the resulting operator
symmetric, hence admissible as a conjugate-gradient preconditioner.

\section{GPU execution}

The hierarchy is built on the host, where the sparse pattern already lives, and
converted to device CSR.  The cycle then runs entirely on the device.  The fine
level is special: rather than storing $\Kmat$ on the device, its
matrix--vector products go through the matrix-free action of
Chapter~\ref{ch:matrixfree}, so the fine matrix is never formed there at all.
Only the coarse levels --- orders of magnitude smaller --- are stored
explicitly.

\section{Lagging and staleness}
\label{sec:amg-lagging}

Hierarchy setup costs seconds at $500$k DOF, far more than a single V-cycle
application, so the hierarchy is \emph{lagged}: built once and reused across
Newton iterations and time steps.  Reuse is safe because the hierarchy is a
preconditioner --- an out-of-date one costs iterations, never accuracy.

Two triggers force a rebuild.  A change in $c_M$ means the operator itself has
changed (a new $\Delta t$).  And a staleness detector watches the
conjugate-gradient iteration count: it latches a baseline on the first solve
after a build and requests a rebuild when a later count exceeds three times
that baseline, subject to a floor so that small counts do not thrash.

On the quasi-static path this detector is the \emph{only} rebuild trigger,
because $c_M = 0$ there and the first test never fires.  That is what makes the
interaction with forcing terms consequential, and it is treated in
Section~\ref{sec:etamax}'s neighbouring discussion at the end of
Chapter~\ref{ch:inexact}.

\section{When multigrid is the right tool}
\label{sec:amg-when}

AMG targets operators whose difficulty is the smooth end of the spectrum:
quasi-static problems and moderate-$\Delta t$ dynamics.  It is not universally
better, and two limits are worth naming.

At small $\Delta t$ in implicit dynamics, $\Keff$ is mass-dominated
(Section~\ref{sec:newmark}) and already well conditioned.  There is little
smooth error for a coarse grid to remove, and measurements on two GPU
architectures confirm the V-cycle does not repay its per-application cost:
Jacobi is faster.

At very large $\Delta t$ on violently dynamic problems, the Newmark predictor
can overshoot into near-inverted configurations whose tangent is indefinite.
Conjugate gradients is then invalid regardless of preconditioner, and explicit
integration is the appropriate method.
